\section{Formal Foundations}\label{sec:foundations}

We formalize decision problems with coordinate structure, sufficiency of coordinate sets, and the optimizer quotient object, drawing on classical decision theory \cite{savage1954foundations, raiffa1961applied}.
The core definitions in this section are substrate-neutral: they specify the decision relation independently of implementation medium. The mechanized physical-lift and claim-transport scaffolding is summarized later in Section~\ref{sec:foundations-mechanized-metadata} so that the mathematical core appears first.

\subsection{Decision Problems with Coordinate Structure}

Formally: this subsection fixes the core tuple, projection, optimizer, and sufficiency/relevance predicates used downstream.

\begin{definition}[Decision Problem]\label{def:decision-problem}
A \emph{decision problem with coordinate structure} is a tuple $\mathcal{D} = (A, X_1, \ldots, X_n, U)$ where:
\begin{itemize}
\item $A$ is a finite set of \emph{actions} (alternatives)
\item $X_1, \ldots, X_n$ are finite \emph{coordinate spaces}
\item $S = X_1 \times \cdots \times X_n$ is the \emph{state space}
\item $U : A \times S \to \mathbb{Q}$ is the \emph{utility function}
\end{itemize}
\end{definition}

\begin{definition}[Projection]\label{def:projection}
For state $s = (s_1, \ldots, s_n) \in S$ and coordinate set $I \subseteq \{1, \ldots, n\}$:
\[
s_I := (s_i)_{i \in I}
\]
is the \emph{projection} of $s$ onto coordinates in $I$.
\end{definition}

\begin{definition}[Optimizer Map]\label{def:optimizer}
For state $s \in S$, the \emph{optimal action set} is:
\[
\Opt(s) := \arg\max_{a \in A} U(a, s) = \{a \in A : U(a,s) = \max_{a' \in A} U(a', s)\}
\]
\end{definition}

\subsection{Sufficiency and Relevance}

\begin{definition}[Sufficient Coordinate Set]\label{def:sufficient}
A coordinate set $I \subseteq \{1, \ldots, n\}$ is \emph{sufficient} for decision problem $\mathcal{D}$ if:
\[
\forall s, s' \in S: \quad s_I = s'_I \implies \Opt(s) = \Opt(s')
\]
Equivalently, the optimal action depends only on coordinates in $I$.
In this paper, this set-valued invariance of the full optimal-action correspondence is the primary decision-relevance target.
\end{definition}

\begin{proposition}[Empty-Set Sufficiency Equals Constant Decision Boundary]\label{prop:empty-sufficient-constant}
For any decision problem,
\[
\text{sufficient}(\emptyset)
\iff
\forall s,s'\in S,\ \Opt(s)=\Opt(s').
\]
Hence, the $I=\emptyset$ query is exactly a universal constancy check of the decision boundary.
\leanmeta{\LH{DP6}.}
\end{proposition}

\begin{proposition}[Insufficiency Equals Counterexample Witness]\label{prop:insufficiency-counterexample}
For any coordinate set $I$,
\[
\neg\text{sufficient}(I)
\iff
\exists s,s'\in S:\ s_I=s'_I\ \wedge\ \Opt(s)\neq\Opt(s').
\]
In particular,
\[
\neg\text{sufficient}(\emptyset)
\iff
\exists s,s'\in S:\ \Opt(s)\neq\Opt(s').
\]
\leanmeta{\LHrng{DP}{7}{8}.}
\end{proposition}
\claimstamp{D\ref{def:sufficient};P\ref{prop:empty-sufficient-constant},\ref{prop:insufficiency-counterexample}}{DM}

\begin{definition}[Deterministic Selector]\label{def:selector}
A \emph{deterministic selector} is a map $\sigma : 2^A \to A$ that returns one action from an optimal-action set.
\end{definition}

\begin{definition}[Selector-Level Sufficiency]\label{def:selector-sufficient}
For fixed selector $\sigma$, a coordinate set $I$ is \emph{selector-sufficient} if
\[
\forall s,s' \in S:\quad s_I=s'_I \implies \sigma(\Opt(s))=\sigma(\Opt(s')).
\]
\end{definition}

\begin{proposition}[Set-Level Target Dominates Selector Targets]\label{prop:set-to-selector}
If $I$ is sufficient in the set-valued sense (Definition~\ref{def:sufficient}), then $I$ is selector-sufficient for every deterministic selector $\sigma$ (Definition~\ref{def:selector-sufficient}).
\leanmeta{\LH{DP5}}
\end{proposition}

\begin{proof}
If $I$ is sufficient, then $s_I=s'_I$ implies $\Opt(s)=\Opt(s')$. Applying any deterministic selector $\sigma$ to equal optimal-action sets yields equal selected actions.
\end{proof}

\begin{definition}[$\varepsilon$-Optimal Set and $\varepsilon$-Sufficiency]\label{def:epsilon-sufficiency}
For $\varepsilon \ge 0$, define
\[
\Opt_\varepsilon(s)
:=
\{a\in A:\forall a'\in A,\ U(a',s)\le U(a,s)+\varepsilon\}.
\]
Coordinate set $I$ is \emph{$\varepsilon$-sufficient} if
\[
\forall s,s'\in S:\ s_I=s'_I \implies \Opt_\varepsilon(s)=\Opt_\varepsilon(s').
\]
\end{definition}

\begin{proposition}[Zero-$\varepsilon$ Reduction]\label{prop:zero-epsilon-reduction}
Set-valued sufficiency is exactly the $\varepsilon=0$ case:
\[
I\text{ sufficient } \iff I\text{ is }0\text{-sufficient}.
\]
\leanmeta{\LH{DP1}, \LH{DP4}}
\end{proposition}

\begin{proposition}[Selector-Level Separation Witness]\label{prop:selector-separation}
The converse of Proposition~\ref{prop:set-to-selector} fails in general: there exists a decision problem, selector, and coordinate set $I$ such that $I$ is selector-sufficient but not sufficient in the full set-valued sense.
\leanmeta{\LH{CC54}.}
\end{proposition}

\begin{definition}[Minimal Sufficient Set]\label{def:minimal-sufficient}
A sufficient set $I$ is \emph{minimal} if no proper subset $I' \subsetneq I$ is sufficient.
\end{definition}

\begin{definition}[Relevant Coordinate]\label{def:relevant}
Coordinate $i$ is \emph{relevant} if there exist states that differ only at coordinate $i$ and induce different optimal action sets:
\[
i \text{ is relevant}
\iff
\exists s,s' \in S:\; \Big(\forall j \neq i,\; s_j = s'_j\Big)\ \wedge\ \Opt(s)\neq\Opt(s').
\]
\end{definition}

\begin{proposition}[Minimal-Set/Relevance Equivalence]\label{prop:minimal-relevant-equiv}
For any minimal sufficient set $I$ and any coordinate $i$:
\[
i \in I \iff i \text{ is relevant}.
\]
Hence every minimal sufficient set is exactly the relevant-coordinate set.
\leanmeta{\LHrng{DP}{2}{3}.}
\end{proposition}

\begin{proof}
The ``only if'' direction follows by minimality: if $i\in I$ were irrelevant, removing $i$ would preserve sufficiency, contradicting minimality. The ``if'' direction follows from sufficiency: every sufficient set must contain each relevant coordinate.
\end{proof}

\begin{proposition}[Structural Rank as Relevant-Support Size]\label{prop:srank-support}
The structural rank $\mathrm{srank}(\mathcal{D})$ is the cardinality of the relevant-coordinate support:
\[
\mathrm{srank}(\mathcal{D})
=
\left|\{i \in \{1,\ldots,n\} : i \text{ is relevant}\}\right|.
\]
Hence $\mathrm{srank}(\mathcal{D})$ is bounded by the ambient coordinate dimension $n$, and $\mathrm{srank}(\mathcal{D}) = 0$ exactly when the decision boundary is constant across coordinates.
\leanmeta{\LHrng{SK}{1}{3}.}
\end{proposition}

\begin{proof}
Proof sketch: by definition, structural rank counts the coordinates that survive the relevance filter. The upper bound follows because at most all coordinates can be relevant. The zero case is equivalent to the relevance support being empty, which is exactly the coordinate-constant boundary case.
\end{proof}

Informally: this is a familiarity anchor for the invariant used throughout the later bridge theorems. In the finite-coordinate setting, $\mathrm{srank}$ is simply the support size of the relevance indicator. Later theorems show that Fisher rank, zero-distortion rate, and physical bit cost recover this same support-size quantity through different frameworks.

\begin{definition}[Exact Relevance Identifiability]\label{def:exact-identifiability}
For a decision problem $\mathcal{D}$ and candidate coordinate set $I$, we say $I$ is \emph{exactly relevance-identifying} if
\[
\forall i \in \{1,\ldots,n\}:\quad i \in I \iff i \text{ is relevant for } \mathcal{D}.
\]
Equivalently, $I$ is exactly relevance-identifying iff $I$ equals the full relevant-coordinate set.
\end{definition}

\begin{example}[Weather Decision]
Consider deciding whether to carry an umbrella:
\begin{itemize}
\item Actions: $A = \{\text{carry}, \text{don't carry}\}$
\item Coordinates: $X_1 = \{\text{rain}, \text{no rain}\}$, $X_2 = \{\text{hot}, \text{cold}\}$, $X_3 = \{\text{Monday}, \ldots, \text{Sunday}\}$
\item Utility: $U(\text{carry}, s) = -1 + 3 \cdot \mathbf{1}[s_1 = \text{rain}]$, $U(\text{don't carry}, s) = -2 \cdot \mathbf{1}[s_1 = \text{rain}]$
\end{itemize}

The minimal sufficient set is $I = \{1\}$ (only rain forecast matters). Coordinates 2 and 3 (temperature, day of week) are irrelevant.
\end{example}

Informally: all later results use this same sufficiency definition.

\subsection{Decision-Quotient Score and Optimizer Quotient Object}

\begin{definition}[Decision Equivalence]\label{def:decision-equiv}
For coordinate set $I$, states $s, s'$ are \emph{$I$-equivalent} (written $s \sim_I s'$) if $s_I = s'_I$.
\end{definition}

\begin{definition}[Decision-Quotient Score]\label{def:decision-quotient}
The \emph{decision-quotient score} for state $s$ under coordinate set $I$ is:
\[
\text{DQ}_I(s) = \frac{|\{a \in A : a \in \Opt(s') \text{ for some } s' \sim_I s\}|}{|A|}
\]
This is a normalized ambiguity score: equivalently, it is the normalized support fraction of the multivalued image of the $I$-equivalence class under $\Opt$. It measures the fraction of actions that are optimal for at least one state consistent with $I$. It is not the quotient object $S/{\sim}$ used later for the optimizer map.
\end{definition}

\begin{proposition}[Sufficiency Characterization]\label{prop:sufficiency-char}
\claimstamp{P\ref{prop:sufficiency-char}}{H=succinct-hard}
Coordinate set $I$ is sufficient if and only if $\text{DQ}_I(s) = |\Opt(s)|/|A|$ for all $s \in S$.
\leanmeta{\LHrng{CC}{62}{63}}
\end{proposition}

\begin{proof}
If $I$ is sufficient, then $s \sim_I s' \implies \Opt(s) = \Opt(s')$, so the set of actions optimal for some $s' \sim_I s$ is exactly $\Opt(s)$.

Conversely, if the condition holds, then for any $s \sim_I s'$, the optimal actions form the same set (since $\text{DQ}_I(s) = \text{DQ}_I(s')$ and both equal the relative size of the common optimal set).
\end{proof}

The quotient object behind the optimizer map has a familiar form in \textbf{Set}. It is useful to state that explicitly before the later abstraction-boundary theorem.

\begin{proposition}[Optimizer Quotient as Coimage/Image Factorization]\label{prop:optimizer-coimage}
For the optimizer map $\Opt : S \to \mathcal{P}(A)$:
\begin{enumerate}
\item the optimizer quotient object is canonically equivalent to $\mathrm{range}(\Opt)$;
\item any surjective abstraction $\phi : S \to T$ that preserves $\Opt$ factors uniquely through this quotient.
\end{enumerate}
Thus, in \textbf{Set}, the optimizer quotient object is the familiar coimage of $\Opt$, canonically identified with its image/range.\cite{maclane1998categories}
\leanmeta{\LH{QT5}, \LH{QT7}.}
\end{proposition}

\begin{proof}
Proof sketch:
\begin{enumerate}
\item The quotient identifies exactly those states with equal optimal-action sets, so its elements correspond bijectively to the values attained by $\Opt$. This gives the canonical equivalence with $\mathrm{range}(\Opt)$.
\item If a surjective abstraction preserves $\Opt$, then the optimizer map already factors through that abstraction. The quotient universal property supplies the mediating map, and surjectivity makes that factorization unique.
\end{enumerate}
\end{proof}

Informally: this proposition is only a familiarity anchor. The object is not exotic in \textbf{Set}; the later novelty is that this quotient governs the complexity and physical-collapse arguments.

The same quotient object has an equally familiar information-theoretic reading: its entropy is just the logarithm of the number of distinct optimizer values actually attained.

\begin{proposition}[Optimizer Quotient Entropy as Image-Size Entropy]\label{prop:optimizer-entropy-image}
Let $\mathrm{numOptClasses}(\mathcal{D})$ denote the number of distinct values attained by $\Opt$, equivalently the cardinality of $\mathrm{range}(\Opt)$. Then
\[
H(\mathcal{D}) = \log_2\bigl(\mathrm{numOptClasses}(\mathcal{D})\bigr)
=
\log_2\bigl(|\mathrm{range}(\Opt)|\bigr).
\]
Thus the entropy of the optimizer quotient object is just the base-2 logarithm of the size of the image of $\Opt$.\leanmeta{\LHrng{IT}{1}{2}, \LH{QT5}.}
\end{proposition}

\begin{proof}
Proof sketch: $\mathrm{numOptClasses}$ is defined as the cardinality of the image of $\Opt$, and Proposition~\ref{prop:optimizer-coimage} identifies the optimizer quotient object with that image in \textbf{Set}. The entropy quantity $H(\mathcal{D})$ is then defined to be $\log_2(\mathrm{numOptClasses})$.
\end{proof}

Informally: this is the familiar reading of the entropy term used later in the bridge theorems. Nothing new is being hidden inside the notation: it is just the logarithm of how many different optimal-action sets the system can realize.

The next theorem concerns the quotient object rather than the scalar score above. Abstraction is a many-to-one quotient operation: it identifies states and thereby removes distinctions. The optimizer quotient object is the coarsest such collapse that still preserves the optimal-action boundary.

Informally: start with any surjective summary map $\phi : S \to T$. There are only two possibilities. Either $\phi$ keeps every optimal-action distinction intact, in which case the universal property forces it to factor through the optimizer quotient object, or $\phi$ merges two states that require different optimal actions, in which case it has erased a decision-relevant distinction. In \textbf{Set}, the quotient object here is the coimage of $\Opt$, canonically equivalent to its image/range \cite{maclane1998categories}. The physical no-collapse layer then says that if one tries to treat that extra erasure as computationally affordable, the collapse assumption itself is what fails.\leanmeta{\LH{QT5}, \LHrng{AB}{1}{4}.}

\begin{theorem}[Decision-Preserving Collapse Boundary]\label{thm:abstraction-boundary}
Let $\phi : S \to T$ be a surjective abstraction of the state space.

\begin{enumerate}
\item Either $\phi$ factors through the optimizer quotient object, or it collapses a decision-relevant distinction, i.e., there exist states $s,s' \in S$ such that $\phi(s)=\phi(s')$ but $\Opt(s)\neq\Opt(s')$.
\item If every such extra collapse is mapped to a physically feasible collapse at the canonical requirement profile, then no such extra collapse can occur. Therefore $\phi$ must factor through the optimizer quotient object.
\end{enumerate}
\leanmeta{\LHrng{AB}{1}{4}.}
\end{theorem}

\begin{proof}
Proof sketch:
\begin{enumerate}
\item If $\phi$ preserves $\Opt$, then the quotient universal property already proved in the previous layer supplies the factorization $\pi = \psi \circ \phi$. In \textbf{Set}, this is exactly the coimage factorization of the optimizer map through its image/range. If $\phi$ does not preserve $\Opt$, then by definition there must exist $s,s'$ with $\phi(s)=\phi(s')$ but $\Opt(s)\neq\Opt(s')$. That witness is exactly the erased decision-relevant distinction.
\item Suppose such an erased distinction is further interpreted as a physically feasible collapse at the canonical requirement profile. The existing physical no-collapse theorem then blocks that collapse map. Therefore the ``extra erasure'' branch is impossible, leaving only factorization through the optimizer quotient object.
\end{enumerate}
\end{proof}

Informally: this is the explicit bridge between state-space abstraction and complexity-class collapse. At the state level, the optimizer quotient object is the maximal lossless collapse. In \textbf{Set}, that means the coimage/image factorization of $\Opt$. At the physical-complexity level, any claim that a coarser collapse remains feasible must pass through the same no-collapse obstruction. The theorem makes the common structure visible instead of leaving it implicit across separate sections.

\subsection{Mechanized Physical-Lift and Claim Metadata}\label{sec:foundations-mechanized-metadata}

The formal system also records the physical-lift and claim-transport scaffolding used later in the paper. Physical interpretation is introduced through explicit encoding and assumption-transfer maps rather than built into the core tuple itself.\claimstamp{D\ref{def:physical-core-encoding};P\ref{prop:physical-claim-transport};T\ref{thm:physical-bridge-bundle}}{AR}
\leanmeta{\LH{CT4}, \LH{CT6}.}

Within that scaffolding, physical-state semantics are typed by realizable-state witnesses, observer classes induce quotient state spaces on a shared physical domain, coupled recurrent circuits support typed YES/NO/ABSTAIN report projection, and timing constraints yield forced-action and lower-bound statements under declared physical parameterizations.
\leanmeta{\LHrng{PS}{1}{4}, \LHrng{OR}{3}{5}, \LH{OR9}, \LHrng{CV}{4}{5}, \LH{CV7}, \LHrng{CV}{11}{13}, \LHrng{CV}{15}{18}, \LH{CF2}, \LH{CF4}, \LHrng{CF}{6}{9}.}

The same metadata layer also packages the assumption-indexed physical-collapse transfer chain, the layered graph witnesses, and the exact-claim admissibility boundary used later in the physics and claim-discipline sections. This material remains explicit, but it is secondary to the mathematical core definitions established above.\claimstamp{T\ref{thm:exact-certified-gap-iff-admissible};P\ref{prop:integrity-resource-bound};C\ref{cor:hardness-raw-only-exact}}{AR}
\leanmeta{\LHrng{PH}{11}{33}, \LH{GN2}, \LHrng{GN}{4}{13}, \LH{CC18}, \LH{CC21}, \LH{CC30}.}

\subsection{Computational Model and Input Encoding}\label{sec:encoding}

Formally: this subsection defines encoding/access regimes, cost measures, and regime-specific complexity interpretation.

We fix the computational model used by the complexity claims.
In physical deployments, these encodings correspond to different access conditions on the same underlying system: explicit tables correspond to fully exposed state/utility structure, succinct encodings correspond to compressed generative descriptions, and query regimes correspond to constrained measurement interfaces.\claimstamp{T\ref{thm:dichotomy};P\ref{prop:query-regime-obstruction},\ref{prop:physical-claim-transport}}{E,S+ETH,Qf,AR}
\leanmeta{\LH{CC16}, \LH{CC50}, \LH{CT6}.}

\paragraph{Succinct encoding (primary for hardness).}
This succinct circuit encoding is the standard representation for decision problems in complexity theory; hardness is stated with respect to the input length of the circuit description \cite{arora2009computational,papadimitriou1994complexity}.
An instance is encoded as:
\begin{itemize}
\item a finite action set $A$ given explicitly,
\item coordinate domains $X_1,\ldots,X_n$ given by their sizes in binary,
\item a Boolean or arithmetic circuit $C_U$ that on input $(a,s)$ outputs $U(a,s)$.
\end{itemize}
The input length is $L = |A| + \sum_i \log |X_i| + |C_U|$. Polynomial time and all complexity classes (\coNP, $\Sigma_2^P$, ETH) are measured in $L$. All hardness results in Section~\ref{sec:hardness} use this encoding.

\paragraph{Explicit-state encoding (used for enumeration algorithms and experiments).}
The utility is given as a full table over $A \times S$. The input length is $L_{\text{exp}} = \Theta(|A||S|)$ (up to the bitlength of utilities). Polynomial time is measured in $L_{\text{exp}}$. Results stated in terms of $|S|$ use this encoding.

\paragraph{Query-access regime (intermediate black-box access).}
The solver is given oracle access to decision information at queried states (e.g., $\Opt(s)$, or $U(a,s)$ with $\Opt(s)$ reconstructed from finitely many value queries). Complexity is measured by oracle-query count (optionally paired with per-query evaluation cost). This separates representational access from full-table availability and from succinct-circuit input length.

Unless explicitly stated otherwise, ``polynomial time'' refers to the succinct encoding.

\begin{remark}[Decision Questions vs Decision Problems]\label{rem:question-vs-problem}
We distinguish between a \emph{decision question} and a \emph{decision problem}. A decision question is the encoding-independent predicate: ``is coordinate set $I$ sufficient for $\mathcal{D}$?'' This is a mathematical relation, fixed by the model contract (C1--C4). A decision problem is a decision question together with a specific encoding of its inputs. The same decision question yields different decision problems under different encodings, and these problems may belong to different complexity classes. Formally, each encoding regime $e$ induces a language $L_{\mathcal{D},e}$; complexity classification is a property of $L_{\mathcal{D},e}$, not of $\mathcal{D}$ alone.
\end{remark}

\begin{definition}[Anchor Query Object]\label{def:anchor-query-object}
An anchor query object is a pair
\[
q=(\mathcal{D},I)
\]
where $\mathcal{D}$ is a decision problem and $I$ is a fixed coordinate set. Its truth predicate is
\[
\mathrm{AnchorTrue}(q)\iff
\exists s_0\in S\ \forall s\in S,\ s_I=s_{0,I}\Rightarrow \Opt(s)=\Opt(s_0).
\]
\end{definition}

\begin{definition}[Pose Operation for Anchor Queries]\label{def:pose-anchor-query}
The pose operation is the constructor
\[
\mathrm{poseAnchor}:\ (\mathcal{D},I)\mapsto q.
\]
Posing is an operation; $q$ is the formal object.
\end{definition}

\begin{proposition}[Posing Produces a Formal Query Object]\label{prop:pose-anchor-object}
For every $(\mathcal{D},I)$, $\mathrm{poseAnchor}(\mathcal{D},I)$ is a well-typed anchor query object with preserved components, and its truth predicate is exactly the anchored $\exists\forall$ condition in Definition~\ref{def:anchor-query-object}.
\claimstamp{D\ref{def:anchor-query-object},\ref{def:pose-anchor-query};T\ref{thm:anchor-sigma2p}}{S}
\leanmeta{\LHrng{AQ}{1}{3}.}
\end{proposition}

\begin{proposition}[Typed Exact Claims on Posed Anchor Queries]\label{prop:posed-anchor-typed-exact}
For posed anchor queries under the typed claim discipline:
\[
\text{Exact admissible}\iff\text{competent on the declared regime},\qquad
\neg\text{competent}\Rightarrow\neg\text{exact admissible}.
\]
Under solver integrity, checker-accepted exact-\texttt{true} verdicts are semantically sound on the posed anchor query object.
\claimstamp{P\ref{prop:typed-claim-admissibility},\ref{prop:integrity-competence-separation};C\ref{cor:no-uncertified-exact-claim}}{AR}
\leanmeta{\LHrng{AQ}{4}{7}.}
\end{proposition}

\begin{corollary}[Certified-Confidence Gate for Posed Anchor Queries]\label{cor:posed-anchor-certified-gate}
Under signal consistency, positive certified confidence implies typed admissibility of the posed anchor-query report.
\claimstamp{D\ref{def:signal-consistency};P\ref{prop:certified-confidence-gate}}{AR}
\leanmeta{\LH{AQ8}.}
\end{corollary}

\begin{proposition}[Under-Resolution Forces Representation Collision]\label{prop:under-resolution-collision}
\claimstamp{P\ref{prop:under-resolution-collision}}{DM}
Let \(\mathrm{repr} : \mathcal{L} \to \mathcal{C}\) be any representation map from logical possibilities to code states. If \(|\mathcal{C}| < |\mathcal{L}|\), then there exist distinct logical states with the same representation:
\[
\exists s,s' \in \mathcal{L},\ s \neq s' \wedge \mathrm{repr}(s)=\mathrm{repr}(s').
\]
This is the finite-cardinality collision statement underlying under-resolved exact-decision claims.
\leanmeta{\LHrng{PI}{6}{7}.}
\end{proposition}

\begin{theorem}[Physical Incompleteness of Instantiation]\label{thm:physical-incompleteness}
\claimstamp{T\ref{thm:physical-incompleteness}}{AR}
Let $\mathcal{P}$ be a finite set of physical states and $\mathcal{L}$ a finite set of logical possibilities, with an instantiation map $\iota : \mathcal{P} \to \mathcal{L}$. If $|\mathcal{P}| < |\mathcal{L}|$, then:
\begin{enumerate}
\item $\iota$ is not surjective;
\item there exists a logical possibility $\ell \in \mathcal{L}$ that is not physically instantiated: $\ell \notin \mathrm{range}(\iota)$.
\end{enumerate}
Under explicit bounds: if $|\mathcal{P}| \le M$ and $|\mathcal{L}| \ge L$ with $M < L$, the same conclusion holds.
\leanmeta{\LHrng{PI}{3}{5}.}
\end{theorem}

\begin{definition}[Regime Simulation]\label{def:regime-simulation}
Fix a decision question $Q$ and two regimes $R_1,R_2$ for that question.
We say $R_1$ \emph{simulates} $R_2$ if there exists a polynomial-time transformer that maps any $R_2$ instance to an $R_1$ instance while preserving the answer to $Q$; in oracle settings, this includes a polynomial-time transcript transducer that preserves yes/no outcomes for induced solvers.
\leanmeta{\LH{CC1}, \LH{CC52}.}
\end{definition}

This paper uses that simulation spine in two explicit forms: restriction-map simulation (subproblem-to-full transfer; Proposition~\ref{prop:query-subproblem-transfer}) and oracle-transducer simulation (batch-to-entry transfer; Proposition~\ref{prop:oracle-lattice-transfer}).

Informally: the same question can be easy or hard depending on representation and access, with explicit transfer requirements.

\subsection{Model Contract and Regime Tags}\label{sec:model-contract}

All theorem statements in this paper are typed by the following model contract:
\begin{itemize}
\item \textbf{C1 (finite actions):} $A$ is finite.
\item \textbf{C2 (finite coordinate domains):} each $X_i$ is finite, so $S=\prod_i X_i$ is finite.
\item \textbf{C3 (evaluable utility):} $U(a,s)$ is computable from the declared input encoding.
\item \textbf{C4 (fixed decision semantics):} optimality is defined by $\Opt(s)=\arg\max_a U(a,s)$.
\end{itemize}

We use short regime tags for applied corollaries:
\begin{itemize}
\item \textbf{[E]} explicit-state encoding,
\item \textbf{[Q]} query-access (oracle) regime,
\item \textbf{[Q\_fin]} finite-state query-access core (state-oracle),
\item \textbf{[S]} succinct encoding,
\item \textbf{[S+ETH]} succinct encoding with ETH,
\item \textbf{[Q\_bool]} mechanized Boolean query submodel,
\item \textbf{[S\_bool]} mechanized Boolean-coordinate submodel.
\end{itemize}

This tagging is a claim-typing convention: each strong statement is attached to the regime where it is proven.

\begin{remark}[Uniform Assumption Discipline]\label{rem:uniform-assumption-discipline}
All theorem-level claims in this paper are conditional on their stated premises and regime tags; no assumption family is given special status.
For each claim, the active assumptions are exactly those appearing in its statement, its regime tag, and its cited mechanized handles.
\end{remark}

\begin{remark}[Modal Convention: Contract-Conditional Necessity]\label{rem:modal-should}
In this paper, normative modal statements denote contract-conditional necessity: once a physical-system model declares an objective/regime/assumption contract, admissible reports and interventions are fixed by typed admissibility and integrity under that contract.
No unconditional external moral prescription is asserted.
\end{remark}

\subsection{Discrete-Time Event Semantics}\label{sec:discrete-time-semantics}

\begin{definition}[Timed Interface State]\label{def:timed-interface-state}
For any interface state space $S$, a timed interface state is a pair
\[
x=(s,t)\in S\times \mathbb{N}.
\]
The time coordinate is the natural-number component $t$.
\end{definition}

\begin{definition}[Decision Tick and Decision Event]\label{def:decision-tick-event}
Fix an interface decision process $(\mathrm{decide},\mathrm{transition})$ with
\[
\mathrm{decide}:S\to A,\qquad \mathrm{transition}:S\times A\to S.
\]
Define one tick by
\[
\mathrm{tick}(s,t)=\bigl(\mathrm{transition}(s,\mathrm{decide}(s)),\,t+1\bigr).
\]
Define $\mathrm{DecisionEvent}(x,y)$ to mean that $y$ is the one-tick successor of $x$ under this rule.
\end{definition}

\begin{proposition}[Time Coordinate Is Discrete]\label{prop:time-discrete}
In Definition~\ref{def:timed-interface-state}, time is discrete:
\[
\forall x\in S\times\mathbb{N},\ \exists k\in\mathbb{N}\ \text{such that}\ x_t=k.
\]
Pointwise time equality is decidable for every $k\in\mathbb{N}$:
\[
x_t=k\ \vee\ x_t\neq k.
\]
\leanmeta{\LHrng{DT}{6}{7}.}
\end{proposition}

\begin{proposition}[Decision Event Is One Unit of Time]\label{prop:decision-unit-time}
Under Definition~\ref{def:decision-tick-event}, for all timed states $x,y$:
\[
\mathrm{DecisionEvent}(x,y)\ \Longleftrightarrow\ y=\mathrm{tick}(x),
\]
and therefore
\[
\mathrm{DecisionEvent}(x,y)\ \Longrightarrow\ y_t=x_t+1.
\]
\leanmeta{\LHrng{DT}{10}{13}.}
\end{proposition}

\begin{proposition}[Run Accounting Identity]\label{prop:run-time-accounting}
Let $\mathrm{run}(n,x)$ be $n$ repeated ticks from $x$, and let $\mathrm{trace}(n,x)$ be the emitted decision sequence. Then:
\[
\mathrm{run}(n,x)_t = x_t + n,\qquad
\mathrm{run}(n,x)_t - x_t = n,\qquad
|\mathrm{trace}(n,x)|=n.
\]
Hence decision count equals elapsed time:
\[
|\mathrm{trace}(n,x)|=\mathrm{run}(n,x)_t-x_t.
\]
\leanmeta{\LHrng{DT}{15}{16}, \LHrng{DT}{18}{19}.}
\end{proposition}

\begin{proposition}[Substrate Realization Preserves Unit-Time Events]\label{prop:substrate-unit-time}
For any substrate model that is consistent with the declared interface transition rule, each one-step substrate evolution realizes one interface decision event and advances interface time by one unit. The time update law is invariant under substrate kind.
\leanmeta{\LHrng{DT}{22}{24}.}
\end{proposition}

\subsection{Physical-System Encoding and Theorem Transport}\label{sec:physical-transport}

Formally: this subsection defines physical-to-core encoding and proves theorem transport via declared assumption maps.

\begin{definition}[Physical-to-Core Encoding]\label{def:physical-core-encoding}
Fix a physical instance class $\mathcal{P}$ and a core decision class $\mathcal{D}$. A \emph{physical-to-core encoding} is a map
\[
E:\mathcal{P}\to\mathcal{D}.
\]
All physical claims below are typed relative to a declared encoding map $E$ and declared assumption transfer map from physical assumptions to core assumptions.
\end{definition}

\begin{definition}[Physical Scope Gate]\label{def:physical-scope-gate}
A claim is in physical scope only if it declares:
\begin{enumerate}
\item a physical state class with observable coordinates,
\item a physical-to-core encoding map \(E:\mathcal{P}\to\mathcal{D}\),
\item calibrated measurement units for reported quantities (for example bits, joules, and \(kT\)-scaled bounds),
\item a declared objective contract and admissible intervention class,
\item an explicit physical-to-core assumption transfer map.
\end{enumerate}
Claims missing any item are untyped for theorem transport in this framework.
\end{definition}

\begin{proposition}[Typed Physical Transport Requirement]\label{prop:typed-physical-transport-requirement}
Direct transfer of core theorems is licensed only through a declared physical scope gate instance (Definition~\ref{def:physical-scope-gate}) and a corresponding physical-to-core encoding (Definition~\ref{def:physical-core-encoding}). Absent that typed physical instance, transport is blocked by construction.
\LH{PS3}\LH{PS4}
\end{proposition}

\begin{proof}
By Definition~\ref{def:physical-scope-gate}, theorem transport is only defined for claims with explicit physical state, objective, intervention, encoding, and assumption-transfer declarations. Therefore untyped targets have no admissible transport map in this framework. Proposition~\ref{prop:physical-claim-transport} then applies only to typed instances.
\end{proof}

This scope gate is a typing rule, not a rhetorical caveat: the paper's transport theorems apply only after explicit encoding and assumption transfer are declared.\claimstamp{D\ref{def:physical-core-encoding},\ref{def:physical-scope-gate};P\ref{prop:typed-physical-transport-requirement},\ref{prop:physical-claim-transport}}{AR}

\begin{proposition}[Heisenberg-Strong Binding Yields Nontrivial Decision Boundary]\label{prop:heisenberg-strong-nontrivial-opt}
\claimstamp{P\ref{prop:heisenberg-strong-nontrivial-opt}}{AR}
Fix a declared noisy physical encoding interface and assume a \emph{Heisenberg-strong binding witness}: one physical instance is compatible with two distinct interface states whose encoded optimal-action sets differ.
Then there exists a core decision problem with non-constant optimizer map.
Equivalently, under this declared physical ambiguity assumption, decision nontriviality is derivable at the core level.
\leanmeta{\LH{HS3}, \LHrng{HS}{5}{6}.}
\end{proposition}

\begin{proposition}[Physical Claim Transport]\label{prop:physical-claim-transport}
Let $E:\mathcal{P}\to\mathcal{D}$ be an encoding. If
\[
\forall d\in\mathcal{D},\ \mathrm{CoreAssm}(d)\Rightarrow \mathrm{CoreClaim}(d)
\]
and
\[
\forall p\in\mathcal{P},\ \mathrm{PhysAssm}(p)\Rightarrow \mathrm{CoreAssm}(E(p)),
\]
then
\[
\forall p\in\mathcal{P},\ \mathrm{PhysAssm}(p)\Rightarrow \mathrm{CoreClaim}(E(p)).
\]
\leanmeta{\LH{CT3}, \LHrng{CT}{5}{6}.}
\end{proposition}

\begin{corollary}[Physical Counterexample Implies Core Failure on Encoded Slice]\label{cor:physical-counterexample-core-failure}
Under the same encoding and assumption-transfer map:
\begin{enumerate}
\item any encoded physical counterexample induces a core counterexample on the encoded slice;
\item if a physical counterexample exists, then the purported core rule is invalid;
\item if the lifted claim holds on all physically admissible encoded instances, no such physical counterexample exists.
\end{enumerate}
\leanmeta{\LHrng{CT}{4}{6}.}
\end{corollary}

\begin{proposition}[Law-Instance Objective Bridge]\label{prop:law-instance-objective-bridge}
For the mechanized law-gap physical instance class, the encoded optimizer claim holds at theorem level and admits no counterexample:
\[
\forall x\in\mathcal{P}_{\mathrm{law}},\ \mathrm{lawGapPhysicalClaim}(x),
\qquad
\neg\exists x\in\mathcal{P}_{\mathrm{law}},\ \neg\,\mathrm{lawGapPhysicalClaim}(x).
\]
\leanmeta{\LHrng{CT}{1}{2}.}
\end{proposition}

\begin{theorem}[Physical Bridge Bundle]\label{thm:physical-bridge-bundle}
In one bundled mechanized statement, the paper's physical bridge layer composes:
\begin{enumerate}
\item law-induced objective semantics (\(\Opt=\) feasible actions for law-gap instances),
\item behavior-preserving configuration reduction equivalence,
\item irreducible required-work lower bound under site-count scaling.
\end{enumerate}
\leanmeta{\LH{CT4}.}
\end{theorem}
\claimstamp{D\ref{def:physical-core-encoding};P\ref{prop:physical-claim-transport},\ref{prop:law-instance-objective-bridge};C\ref{cor:physical-counterexample-core-failure};T\ref{thm:physical-bridge-bundle}}{AR}

\begin{theorem}[Declared Regime Coverage for the Static Class]\label{thm:regime-coverage}
Let
\[
\mathcal{R}_{\mathrm{static}} =
\{\text{[E]},\ \text{[S+ETH]},\ \text{[Q\_fin:Opt]},\ \text{[Q\_bool:value-entry]},\ \text{[Q\_bool:state-batch]}\}.
\]
For each declared regime in $\mathcal{R}_{\mathrm{static}}$, the paper has a theorem-level mechanized core claim. Equivalently, regime typing is complete over the declared static family.
\leanmeta{\LH{CC14}, \LH{CC51}.}
\end{theorem}

\paragraph{Scope Lattice (typed classes and transfer boundary).}
\begin{center}
\small
\setlength{\tabcolsep}{4pt}
\begin{tabularx}{\linewidth}{@{}>{\raggedright\arraybackslash}p{0.37\linewidth}>{\raggedright\arraybackslash}p{0.18\linewidth}>{\raggedright\arraybackslash}X@{}}
\toprule
\textbf{Layer} & \textbf{Transfer Status} & \textbf{Lean handles (representative)} \\
\midrule
Static sufficiency class (C1--C4; declared regimes) & Internal landscape complete & \LH{CC76}\\\LH{CC51} \\
Bridge-admissible adjacent class (one-step deterministic) & Transfer licensed & \LH{CC42}\\\LH{CC10} \\
Non-admissible represented adjacent classes (horizon, stochastic, transition-coupled) & Transfer blocked by witness & \LH{CC9}\\\LH{CC8} \\
\bottomrule
\end{tabularx}
\end{center}

Informally: physical conclusions apply when the stated physical assumptions line up with the core theorem assumptions.

\subsection{Adjacent Objective Regimes and Bridge}

\begin{definition}[Adjacent Sequential Objective Regime]\label{def:adjacent-sequential-regime}
An adjacent sequential objective instance consists of:
\begin{itemize}
\item finite action set $A$,
\item finite coordinate state space $S = X_1 \times \cdots \times X_n$,
\item horizon $T \in \mathbb{N}_{\ge 1}$ and history-dependent control-law class,
\item reward process $r_t$ and objective functional $J(\pi)$ (e.g., cumulative reward or regret).
\end{itemize}
\end{definition}

\begin{proposition}[One-Step Deterministic Bridge]\label{prop:one-step-bridge}
Consider an instance of Definition~\ref{def:adjacent-sequential-regime} satisfying:
\begin{enumerate}
\item $T=1$,
\item deterministic rewards $r_1(a,s)=U(a,s)$ for some evaluable $U$,
\item objective $J(\pi)=U(\pi(s),s)$ (single-step maximization),
\item no post-decision state update relevant to the objective.
\end{enumerate}
Then the induced optimization problem is exactly the static decision problem of Definition~\ref{def:decision-problem}, and coordinate sufficiency in the sequential formulation is equivalent to Definition~\ref{def:sufficient}.
\leanmeta{\LH{CC42}}
\end{proposition}

\begin{proof}
Under (1)--(3), optimizing $J$ at state $s$ is identical to choosing an action in $\arg\max_{a\in A} U(a,s)=\Opt(s)$. Condition (4) removes any dependence on future transition effects. Therefore the optimal-policy relation in the adjacent formulation coincides pointwise with $\Opt$ from Definition~\ref{def:optimizer}, and the invariance condition ``same projection implies same optimal choice set'' is exactly Definition~\ref{def:sufficient}.
\end{proof}

\begin{proposition}[Bridge Transfer Rule]\label{prop:bridge-transfer-scope}
Under conditions (1)--(4), any sufficiency statement formulated over Definition~\ref{def:sufficient} is equivalent between the adjacent sequential formulation and the static model.
\leanmeta{\LH{CC42}}
\end{proposition}

\begin{proof}
By Proposition~\ref{prop:one-step-bridge}, the adjacent sequential objective induces exactly the same sufficiency predicate as Definition~\ref{def:sufficient} when (1)--(4) hold. Equivalence of any sufficiency statement then follows by substitution.
\end{proof}

If any bridge condition fails, direct transfer from this paper's static complexity theorems is not licensed by this rule.

\begin{remark}[Extension Boundary]
Beyond Proposition~\ref{prop:one-step-bridge} (multi-step horizon, stochastic transitions/rewards, or regret objectives), the governing complexity objects change. Those regimes are natural extensions, but they are distinct formal classes from the static sufficiency class analyzed in this paper.
\end{remark}

\begin{proposition}[Horizon-$>1$ Bridge-Failure Witness]\label{prop:bridge-failure-horizon}
Dropping the one-step condition can break sufficiency transfer: there exists a two-step objective where sufficiency in the immediate-utility static projection does not match sufficiency in the two-step objective.
\leanmeta{\LHrng{CC}{25}{26}.}
\end{proposition}

\begin{proposition}[Stochastic-Criterion Bridge-Failure Witness]\label{prop:bridge-failure-stochastic}
If optimization is performed against a stochastic criterion not equal to deterministic utility maximization, bridge transfer to Definition~\ref{def:decision-problem} can fail even when an expected-utility projection is available.
\leanmeta{\LH{CC57}.}
\end{proposition}

\begin{proposition}[Transition-Coupled Bridge-Failure Witness]\label{prop:bridge-failure-transition}
If post-decision transition effects are objective-relevant, bridge transfer to the static one-step class can fail.
\leanmeta{\LH{CC73}.}
\end{proposition}

\begin{theorem}[Exact Bridge Boundary in the Represented Adjacent Family]\label{thm:bridge-boundary-represented}
Within the represented adjacent family
\[
\{\text{one-step deterministic},\ \text{horizon-extended},\ \text{stochastic-criterion},\ \text{transition-coupled}\},
\]
direct transfer of static sufficiency claims is licensed if and only if the class is one-step deterministic. Each represented non-one-step class has an explicit mechanized bridge-failure witness.
\leanmeta{\LHrng{CC}{8}{10}.}
\end{theorem}

\begin{definition}[System Snapshot vs System Process Typing]\label{def:system-snapshot-process}
Within the represented adjacent family, we type:
\begin{itemize}
\item \emph{system snapshot}: fixed-parameter inference view (mapped to one-step deterministic class),
\item \emph{system process}: online/dynamical update views (horizon-extended, stochastic-criterion, or transition-coupled classes).
\end{itemize}
\end{definition}

\begin{proposition}[Snapshot/Process Transfer Typing]\label{prop:snapshot-process-typing}
In the represented adjacent family, direct transfer of static sufficiency claims is licensed if and only if the system is typed as a \emph{system snapshot}. Every process-typed represented class has an explicit mechanized bridge-failure witness.
\leanmeta{\LH{CC3}, \LH{CC48}, \LH{CC56}.}
\end{proposition}

Operationally: a fixed-parameter physical controller at runtime is a system snapshot for this typing; once parameters or transition-relevant internals are updated online, the object is process-typed and transfer from static sufficiency theorems is blocked unless the one-step bridge conditions are re-established.\claimstamp{P\ref{prop:snapshot-process-typing}}{RA}

\section{Interpretive Foundations: Hardness and Solver Claims}\label{sec:interpretive-foundations}

The claims in later applied sections are theorem-indexed consequences of this section and Sections~\ref{sec:hardness}--\ref{sec:tractable}.

\subsection{Structural Complexity vs Representational Hardness}

Formally: this subsection separates class placement from regime-indexed access cost and states typed completeness for the declared static family.

\begin{definition}[Structural Complexity]\label{def:structural-complexity}
For a fixed decision question (e.g., ``is $I$ sufficient for $\mathcal{D}$?''; see Remark~\ref{rem:question-vs-problem}), \emph{structural complexity} means its placement in standard complexity classes within the formal model (coNP, $\Sigma_2^P$, etc.), as established by class-membership arguments and reductions.
\end{definition}

\begin{definition}[Representational Hardness]\label{def:representational-hardness}
For a fixed decision question and an encoding regime $E$ (Section~\ref{sec:encoding}), \emph{representational hardness} is the worst-case computational cost incurred by solvers whose input access is restricted to $E$.
\end{definition}

In the mechanized architecture model, this is made explicit as a required-work observable and intrinsic lower bound: `requiredWork` is total realized work, `hardnessLowerBound` is the intrinsic floor, and `hardness\_is\_irreducible\_required\_work` proves the lower-bound relation for all nonzero use-site counts; `requiredWork\_eq\_affine\_in\_sites` and `workGrowthDegree` make the site-count growth law explicit.\leanmeta{\LHrng{HD}{17}{18}, \LHrng{HD}{22}{23}, \LHrng{HD}{27}{28}.}

\begin{remark}[Interpretation Contract]
This paper keeps the decision question fixed and varies the encoding regime explicitly. Thus, later separations are read as changes in representational hardness under fixed structural complexity, not as changes to the underlying sufficiency semantics.
Accordingly, ``complete landscape'' means complete for this static sufficiency class under the declared regimes; adjacent objective classes are distinct typed objects, not unproved remainder cases of the same class.
\end{remark}

\begin{theorem}[Typed Completeness for the Static Sufficiency Class]\label{thm:typed-completeness-static}
For the static sufficiency class fixed by C1--C4 and declared regime family $\mathcal{R}_{\mathrm{static}}$, the paper provides:
\begin{enumerate}
\item conditional class-label closures for SUFFICIENCY-CHECK, MINIMUM-SUFFICIENT-SET, and ANCHOR-SUFFICIENCY;
\item a conditional explicit/succinct dichotomy closure;
\item a regime-indexed mechanized core claim for every declared regime; and
\item declared-family exhaustiveness in the typed regime algebra.
\end{enumerate}
\leanmeta{\LH{CC76}.}
\end{theorem}

Informally: the question stays fixed while representation changes cost, and the declared static regime set is fully covered.

\subsection{Solver Integrity and Regime Competence}

Formally: this subsection defines integrity, competence, typed report/evidence admissibility, and resource-bounded exact-claim limits.

To keep practical corollaries type-safe, we separate \emph{integrity} (what a solver is allowed to assert) from \emph{competence} (what it can cover under a declared regime), following the certifying-algorithms schema \cite{mcconnell2010certifying,necula1997proof}.
The induced-relation view used below is standard in complexity/computability presentations: algorithms are analyzed as machines deciding languages or computing functions/relations over encodings \cite{papadimitriou1994complexity,arora2009computational,forster2019verified}.

\begin{definition}[Certifying Solver]\label{def:certifying-solver}
Fix a decision question $\mathcal{R} \subseteq \mathcal{X}\times\mathcal{Y}$ and an encoding regime $E$ over $\mathcal{X}$. A \emph{certifying solver} is a pair $(Q,V)$ where:
\begin{itemize}
\item $Q$ maps each input $x\in\mathcal{X}$ to either $\mathsf{ABSTAIN}$ or a candidate pair $(y,w)$,
\item $V$ is a polynomial-time checker (in $|{\rm enc}_E(x)|$) with output in $\{0,1\}$.
\end{itemize}
\end{definition}

\begin{definition}[Induced Solver Relation]\label{def:induced-solver-relation}
For any deterministic program $M$ computing a partial map
$f_M:\mathcal{X}\rightharpoonup\mathcal{Y}$ on encoded inputs, define
\[
\mathcal{R}_M
:=
\{(x,y)\in\mathcal{X}\times\mathcal{Y} : f_M(x)\downarrow \ \wedge\ y=f_M(x)\}.
\]
\end{definition}

\begin{proposition}[Universality of Solver Framing]\label{prop:universal-solver-framing}
Every deterministic computer program that computes a (possibly partial) map on encoded inputs can be framed as a solver for its induced relation (Definition~\ref{def:induced-solver-relation}).
In this formal sense, all computers are solvers of specific problems.
\leanmeta{\LH{CC77}.}
\end{proposition}

\begin{proof}
Immediate from Definition~\ref{def:induced-solver-relation}: each program induces a relation given by its graph on the domain where it halts.
\end{proof}

\begin{definition}[Solver Integrity]\label{def:solver-integrity}
A certifying solver $(Q,V)$ has \emph{integrity} for relation $\mathcal{R}$ if:
\begin{itemize}
\item (assertion soundness) $Q(x)=(y,w)\implies V(x,y,w)=1$,
\item (checker soundness) $V(x,y,w)=1\implies (x,y)\in\mathcal{R}$.
\end{itemize}
The output $\mathsf{ABSTAIN}$ (equivalently, $\mathsf{UNKNOWN}$) is first-class and carries no assertion about membership in $\mathcal{R}$.
\end{definition}

\begin{definition}[Competence Under a Regime]\label{def:competence-regime}
Fix a regime $\Gamma=(\mathcal{X}_\Gamma,E_\Gamma,\mathcal{C}_\Gamma)$ with instance family $\mathcal{X}_\Gamma\subseteq\mathcal{X}$, encoding assumptions $E_\Gamma$, and resource bound $\mathcal{C}_\Gamma$. A certifying solver $(Q,V)$ is \emph{competent} on $\Gamma$ for relation $\mathcal{R}$ if:
\begin{itemize}
\item it has integrity for $\mathcal{R}$ (Definition~\ref{def:solver-integrity}),
\item (coverage) $\forall x\in\mathcal{X}_\Gamma,\;Q(x)\neq\mathsf{ABSTAIN}$,
\item (resource bound) runtime$_Q(x)\le \mathcal{C}_\Gamma(|{\rm enc}_{E_\Gamma}(x)|)$ for all $x\in\mathcal{X}_\Gamma$.
\end{itemize}
\end{definition}

\begin{corollary}[Universality of Integrity Applicability]\label{cor:integrity-universal}
Let $M$ be any deterministic program, viewed as a certifying solver pair $(Q,V)$ for some externally fixed target relation $\mathcal{R}$.
Then Definition~\ref{def:solver-integrity} applies unchanged.
Thus epistemic integrity is architecture-universal over computing systems once they are cast as solvers for declared tasks.
\leanmeta{\LH{CC31}.}
\end{corollary}

\begin{proof}
Definition~\ref{def:solver-integrity} quantifies over a relation $\mathcal{R}$ and a certifying pair $(Q,V)$; it does not assume any special implementation substrate.
\end{proof}
\claimstamp{P\ref{prop:universal-solver-framing};C\ref{cor:integrity-universal}}{TR}

\begin{remark}[External-Task Non-Vacuity]\label{rem:external-task-nonvacuity}
If the target relation is chosen post hoc from the program's own behavior (for example, $\mathcal{R}=\mathcal{R}_M$ from Definition~\ref{def:induced-solver-relation}), integrity can become tautological and competence claims can be vacuous.
Non-vacuous competence claims therefore require the target relation, encoding regime, and resource bound to be declared independently of observed outputs.
\claimstamp{D\ref{def:solver-integrity},\ref{def:competence-regime};P\ref{prop:integrity-competence-separation},\ref{prop:integrity-resource-bound}}{AR}
\end{remark}

\begin{definition}[$\varepsilon$-Objective Relation Family]\label{def:epsilon-objective-family}
An \emph{$\varepsilon$-objective relation family} is a map
\[
\varepsilon \mapsto \mathcal{R}_\varepsilon \subseteq \mathcal{X}\times\mathcal{Y}
\]
that assigns a certified target relation to each tolerance level $\varepsilon\ge 0$.
The exact target is $\mathcal{R}_0$.
\end{definition}

\begin{definition}[$\varepsilon$-Competence Under a Regime]\label{def:epsilon-competence-regime}
Fix a relation family $(\mathcal{R}_\varepsilon)_{\varepsilon\ge 0}$ and regime $\Gamma$.
A certifying solver $(Q,V)$ is \emph{$\varepsilon$-competent} on $\Gamma$ if it is competent on $\Gamma$ for relation $\mathcal{R}_\varepsilon$ in the sense of Definition~\ref{def:competence-regime}.
\end{definition}

\begin{proposition}[Zero-$\varepsilon$ Competence Reduction]\label{prop:zero-epsilon-competence}
If $\mathcal{R}_0=\mathcal{R}$, then
\[
\text{$\varepsilon$-competent at }\varepsilon=0
\iff
\text{competent for exact relation }\mathcal{R}.
\]
Moreover, $\varepsilon$-competence implies integrity for the corresponding $\mathcal{R}_\varepsilon$.
\leanmeta{\LH{IC19}, \LH{IC33}.}
\end{proposition}

\begin{definition}[Typed Claim Report]\label{def:typed-claim-report}
For a declared objective family and regime, a solver-side report is one of:
\begin{itemize}
\item $\mathsf{ABSTAIN}$,
\item $\mathsf{EXACT}$ (claiming exact competence/certification),
\item $\mathsf{EPSILON}(\varepsilon)$ (claiming $\varepsilon$-competence/certification).
\end{itemize}
\end{definition}

\begin{proposition}[Typed Claim Admissibility Discipline]\label{prop:typed-claim-admissibility}
Under Definitions~\ref{def:competence-regime}, \ref{def:epsilon-competence-regime}, and \ref{def:typed-claim-report}:
\begin{itemize}
\item $\mathsf{ABSTAIN}$ is always admissible;
\item $\mathsf{EXACT}$ is admissible iff exact competence holds;
\item $\mathsf{EPSILON}(\varepsilon)$ is admissible iff $\varepsilon$-competence holds.
\end{itemize}
Hence claim type and certificate type must match (no cross-typing of uncertified assertions as certified guarantees).
\leanmeta{\LH{CC75}.}
\end{proposition}

\begin{definition}[Evidence Object for a Typed Report]\label{def:evidence-object}
For declared $(\mathcal{R},(\mathcal{R}_\varepsilon)_{\varepsilon\ge 0},\Gamma,Q)$ and report type $r\in\{\mathsf{ABSTAIN},\mathsf{EXACT},\mathsf{EPSILON}(\varepsilon)\}$,
\emph{evidence} is a first-class witness object:
\begin{itemize}
\item for $\mathsf{ABSTAIN}$: trivial witness;
\item for $\mathsf{EXACT}$: exact-competence witness;
\item for $\mathsf{EPSILON}(\varepsilon)$: $\varepsilon$-competence witness.
\end{itemize}
\leanmeta{\LH{IC2}.}
\end{definition}

\begin{proposition}[Evidence--Admissibility Equivalence]\label{prop:evidence-admissibility-equivalence}
For any declared contract and report type $r$:
\[
\text{Evidence object for }r\ \text{exists}
\iff
\text{Claim }r\ \text{is admissible}.
\]
\leanmeta{\LH{IC17}, \LHrng{IC}{20}{21}.}
\end{proposition}
\claimstamp{D\ref{def:evidence-object};P\ref{prop:typed-claim-admissibility},\ref{prop:evidence-admissibility-equivalence}}{AR}

\paragraph{Security-game reading (derived).}
The typed-report layer can be read as a standard claim-verification game under a declared contract:
a reporting algorithm emits a report token $r$ and optional evidence, and a checker accepts if and only if the evidence is valid for $r$.
Within this game view, existing results already provide the core security properties:
\emph{completeness} for admissible report classes (Proposition~\ref{prop:evidence-admissibility-equivalence}),
\emph{soundness} against overclaiming (Propositions~\ref{prop:typed-claim-admissibility}, \ref{prop:exact-requires-evidence}),
and \emph{evidence-gated confidence} (Propositions~\ref{prop:certified-confidence-gate}, \ref{prop:no-evidence-zero-certified}).
So integrity is enforced as verifiable claim discipline, not as unchecked payload declaration.
\claimstamp{P\ref{prop:typed-claim-admissibility},\ref{prop:evidence-admissibility-equivalence},\ref{prop:exact-requires-evidence},\ref{prop:certified-confidence-gate},\ref{prop:no-evidence-zero-certified}}{AR}

\begin{definition}[Signaled Typed Report (Guess + Confidence)]\label{def:signaled-typed-report}
A \emph{signaled typed report} augments the typed report token
$r\in\{\mathsf{ABSTAIN},\mathsf{EXACT},\mathsf{EPSILON}(\varepsilon)\}$ with:
\begin{itemize}
\item an optional guess payload $g$,
\item self-reflected confidence $p_{\mathrm{self}}\in[0,1]$,
\item certified confidence $p_{\mathrm{cert}}\in[0,1]$.
\item scalar execution witness $t_{\mathrm{run}}\in\mathbb{N}$ (steps actually run),
\item optional declared bound $B\in\mathbb{N}$.
\end{itemize}
Here $p_{\mathrm{self}}$ is a self-signal, while $p_{\mathrm{cert}}$ is an evidence-gated signal under the declared contract.
\leanmeta{\LH{IC34}.}
\end{definition}

\begin{definition}[Signal Consistency Discipline]\label{def:signal-consistency}
For a signaled typed report $(r,g,p_{\mathrm{self}},p_{\mathrm{cert}},t_{\mathrm{run}},B)$, require:
\[
p_{\mathrm{cert}} > 0 \Rightarrow \exists\ \text{evidence object for }r.
\]
\end{definition}

\begin{proposition}[Exact Claims Require Exact Evidence]\label{prop:exact-requires-evidence}
In the typed claim discipline,
\[
\text{ClaimAdmissible}(\mathsf{EXACT})
\iff
\exists\ \text{exact evidence object}.
\]
Equivalently, admissible exact claims are exactly those with an evidence witness.
\leanmeta{\LHrng{IC}{35}{36}.}
\end{proposition}

\begin{definition}[Bound-Dependent Completion Fraction]\label{def:completion-fraction}
Completion-fraction semantics are defined only when a positive declared bound exists and covers observed runtime:
\[
\mathrm{CompletionFractionDefined}
\iff
\exists b>0:\ B=b\ \wedge\ t_{\mathrm{run}}\le b.
\]
\leanmeta{\LH{IC37}.}
\end{definition}

\begin{proposition}[Certified-Confidence Gate]\label{prop:certified-confidence-gate}
Under Definition~\ref{def:signal-consistency}, positive certified confidence implies typed admissibility:
\[
p_{\mathrm{cert}} > 0 \Rightarrow \text{ClaimAdmissible}(r).
\]
Thus certified confidence is not a free scalar; it is evidence-gated by the same typed discipline as the report itself.
\leanmeta{\LHrng{IC}{38}{39}.}
\end{proposition}

\begin{proposition}[No Evidence Forces Zero Certified Confidence]\label{prop:no-evidence-zero-certified}
Under Definition~\ref{def:signal-consistency}, absence of evidence for the reported type forces:
\[
\neg\exists\ \text{evidence object for }r
\Rightarrow
p_{\mathrm{cert}}=0.
\]
\leanmeta{\LH{IC40}.}
\end{proposition}

\begin{corollary}[Exact-No-Competence Zero-Certified Constraint]\label{cor:exact-no-competence-zero-certified}
For exact reports, if exact competence is unavailable on the declared regime/objective, then any signal-consistent report must satisfy:
\[
p_{\mathrm{cert}} = 0.
\]
\leanmeta{\LH{IC41}.}
\end{corollary}

\begin{proposition}[Scalar Step Witness Is Always Defined and Falsifiable]\label{prop:steps-run-scalar}
For every signaled report, the execution witness is an exact natural number and any equality claim about it is decidable:
\[
\exists k\in\mathbb{N}:\ t_{\mathrm{run}}=k,
\qquad
\forall k\in\mathbb{N},\ (t_{\mathrm{run}}=k)\ \vee\ (t_{\mathrm{run}}\neq k).
\]
\leanmeta{\LHrng{IC}{42}{43}.}
\end{proposition}

\begin{proposition}[No Completion Fraction Without Declared Bound]\label{prop:no-fraction-without-bound}
If no declared bound is provided ($B=\varnothing$), completion-fraction semantics are undefined:
\[
B=\varnothing \Rightarrow \neg\,\mathrm{CompletionFractionDefined}.
\]
\leanmeta{\LH{IC44}.}
\end{proposition}

\begin{proposition}[Completion Fraction Under Explicit Bound]\label{prop:fraction-defined-under-bound}
If a positive declared bound is provided and observed runtime is within bound, completion-fraction semantics are defined:
\[
B=b,\ b>0,\ t_{\mathrm{run}}\le b
\Rightarrow
\mathrm{CompletionFractionDefined}.
\]
\leanmeta{\LH{IC45}.}
\end{proposition}

\begin{proposition}[Abstain Can Carry Guess and Self-Confidence]\label{prop:abstain-guess-self-signal}
For any optional guess payload $g$ and any self-reflected confidence $p_{\mathrm{self}}\in[0,1]$, there exists a signal-consistent abstain report
\[
(\mathsf{ABSTAIN}, g, p_{\mathrm{self}}, 0, 0, \varnothing).
\]
Hence the framework is non-binary at the report-signal layer: abstention can carry a tentative answer and self-reflection without upgrading to certified exactness.
\leanmeta{\LH{IC46}.}
\end{proposition}

\begin{proposition}[Self-Reflected Confidence Is Not Certification]\label{prop:self-confidence-not-certification}
Self-reflected confidence alone does not certify exactness: if exact competence is unavailable, there exist exact-labeled reports with arbitrary $p_{\mathrm{self}}$ that remain inadmissible under zero certified confidence.
\leanmeta{\LH{IC47}.}
\end{proposition}
\claimstamp{D\ref{def:signaled-typed-report},\ref{def:signal-consistency},\ref{def:completion-fraction};P\ref{prop:exact-requires-evidence},\ref{prop:certified-confidence-gate},\ref{prop:no-evidence-zero-certified},\ref{prop:steps-run-scalar},\ref{prop:no-fraction-without-bound},\ref{prop:fraction-defined-under-bound},\ref{prop:abstain-guess-self-signal},\ref{prop:self-confidence-not-certification};C\ref{cor:exact-no-competence-zero-certified}}{AR}

\begin{definition}[Declared Budget Slice]\label{def:declared-budget-slice}
For a declared regime $\Gamma$ over state objects and horizon/budget cut $H\in\mathbb{N}$, define the in-scope slice
\[
\mathcal{S}_{\Gamma,H}
:=
\{s:\ s\in\Gamma.\mathrm{inScope}\ \wedge\ \mathrm{encLen}_\Gamma(s)\le H\}.
\]
\leanmeta{\LH{CC13}.}
\end{definition}

\begin{proposition}[Bounded-Slice Irrelevance of a Meta-Coordinate]\label{prop:bounded-slice-meta-irrelevance}
Fix coordinate $i_\infty$ and declared slice $\mathcal{S}_{\Gamma,H}$ from Definition~\ref{def:declared-budget-slice}.
If optimizer sets are invariant to $i_\infty$ on that slice, i.e.,
\[
\forall s,s'\in\mathcal{S}_{\Gamma,H},\ 
\Big(\forall j\neq i_\infty,\ s_j=s'_j\Big)
\Rightarrow
\Opt(s)=\Opt(s'),
\]
then $i_\infty$ is irrelevant on $\mathcal{S}_{\Gamma,H}$, and hence not relevant for the declared task slice.
\leanmeta{\LHrng{CC}{32}{33}.}
\end{proposition}
\claimstamp{D\ref{def:declared-budget-slice};P\ref{prop:bounded-slice-meta-irrelevance}}{AR}

\begin{definition}[Goal Class]\label{def:goal-class}
A \emph{goal class} is a non-empty set of admissible evaluators over a fixed action/state structure:
\[
\mathcal{G}=(\mathcal{U}_{\mathcal{G}},A,S),\qquad
\varnothing\neq\mathcal{U}_{\mathcal{G}}\subseteq\{U:A\times S\to\mathbb{Q}\}.
\]
The solver need not know which $U\in\mathcal{U}_{\mathcal{G}}$ is active; it only knows class-membership constraints.
\leanmeta{\LH{IV1}.}
\end{definition}

\begin{definition}[Interior Pareto Dominance]\label{def:interior-pareto-dominance}
For a goal class $\mathcal{G}$ and score model $\sigma:S\times\{1,\ldots,n\}\to\mathbb{Q}$, let $J_{\mathcal{G}}$ be the class-tautological coordinate set.
State $s$ \emph{interior-Pareto-dominates} $s'$ on $J_{\mathcal{G}}$ when:
\[
\forall j\in J_{\mathcal{G}},\ \sigma(s',j)\le \sigma(s,j)
\quad\text{and}\quad
\exists j\in J_{\mathcal{G}},\ \sigma(s',j)<\sigma(s,j).
\]
\leanmeta{\LHrng{IV}{2}{3}.}
\end{definition}

\begin{proposition}[Interior Dominance Implies Universal Non-Rejection]\label{prop:interior-universal-non-rejection}
If a state $s$ interior-Pareto-dominates $s'$ on a checked coordinate set $J$, and every admissible evaluator in the goal class is monotone on $J$, then no admissible evaluator strictly prefers $s'$ over $s$ on that checked slice:
\[
\forall U\in\mathcal{U}_{\mathcal{G}},\ \forall a\in A,\ U(a,s')\le U(a,s).
\]
\leanmeta{\LH{IV5}.}
\end{proposition}

\begin{proposition}[Interior Verification Tractability]\label{prop:interior-verification-tractable}
\claimstamp{P\ref{prop:interior-verification-tractable}}{H=tractable-structured}
If membership in the checked coordinate set is decidable and interior dominance is decidable, then interior verification yields a polynomial-time checker with exact acceptance criterion:
\[
\exists\ \mathrm{verify}:S\times S\to\{0,1\},\quad
\mathrm{verify}(s,s')=1\iff \text{interior-dominance}(s,s').
\]
\leanmeta{\LH{IV4}, \LH{IV7}.}
\end{proposition}

\begin{proposition}[Interior Dominance Is One-Sided]\label{prop:interior-one-sidedness}
Interior dominance is one-sided and does not imply full coordinate sufficiency:
a counterexample pair outside the checked slice can still violate global optimizer equivalence.
\leanmeta{\LH{IV6}, \LH{IV8}.}
\end{proposition}

\begin{corollary}[Singleton Interior Certificate]\label{cor:interior-singleton-certificate}
\claimstamp{C\ref{cor:interior-singleton-certificate}}{H=tractable-structured}
For a singleton checked coordinate $i$, strict improvement on $i$ with class-monotonicity on $\{i\}$ is already a valid interior certificate of non-rejection.
\leanmeta{\LH{IV9}.}
\end{corollary}
\claimstamp{D\ref{def:goal-class},\ref{def:interior-pareto-dominance};P\ref{prop:interior-universal-non-rejection},\ref{prop:interior-verification-tractable},\ref{prop:interior-one-sidedness};C\ref{cor:interior-singleton-certificate}}{AR}

\begin{definition}[RLFF Objective (Evidence-Gated Report Reward)]\label{def:rlff-objective}
Given report type $r\in\{\mathsf{ABSTAIN},\mathsf{EXACT},\mathsf{EPSILON}(\varepsilon)\}$ and declared contract,
RLFF assigns base reward by report type when evidence exists and applies an explicit inadmissibility penalty otherwise:
\[
\mathrm{Reward}(r)=
\begin{cases}
\mathrm{Base}(r), & \text{if evidence for }r\text{ exists},\\
-\mathrm{Penalty}, & \text{otherwise}.
\end{cases}
\]
\leanmeta{\LH{IC4}, \LHrng{IC}{28}{29}.}
\end{definition}

\begin{proposition}[RLFF Maximizer Is Admissible]\label{prop:rlff-maximizer-admissible}
If $\mathsf{ABSTAIN}$ reward strictly exceeds the inadmissibility branch, then any global RLFF maximizer must be evidence-backed and therefore admissible under the typed-claim discipline.
\leanmeta{\LHrng{IC}{31}{32}.}
\end{proposition}

\begin{corollary}[RLFF Abstain Preference Without Certificates]\label{cor:rlff-abstain-no-certs}
If $\mathsf{EXACT}$ and $\mathsf{EPSILON}(\varepsilon)$ have no evidence objects while $\mathsf{ABSTAIN}$ beats the inadmissibility branch, then RLFF strictly prefers $\mathsf{ABSTAIN}$ to both report types.
\leanmeta{\LH{IC30}, \LH{IC40}.}
\end{corollary}
\claimstamp{D\ref{def:rlff-objective};P\ref{prop:rlff-maximizer-admissible};C\ref{cor:rlff-abstain-no-certs}}{AR}

\begin{definition}[Certainty Inflation]\label{def:certainty-inflation}
For a typed report $r$, certainty inflation is the state of emitting $r$ without a matching evidence object.
\leanmeta{\LH{IC1}, \LH{IC3}.}
\end{definition}

\begin{proposition}[Certainty Inflation Equals Report Inadmissibility]\label{prop:certainty-inflation-iff-inadmissible}
For every typed report $r$:
\[
\text{CertaintyInflation}(r)
\iff
\neg\ \text{ClaimAdmissible}(r).
\]
For $r=\mathsf{EXACT}$ this is equivalently failure of exact competence.
\leanmeta{\LH{IC8}, \LH{IC22}.}
\end{proposition}

\begin{corollary}[Hardness-Blocked Regimes Induce Exact Certainty Inflation]\label{cor:hardness-exact-certainty-inflation}
Under the declared hardness/non-collapse premises that block universal exact admissibility, every exact report is certainty-inflated.
\leanmeta{\LH{CC19}.}
\end{corollary}
\claimstamp{D\ref{def:certainty-inflation};P\ref{prop:certainty-inflation-iff-inadmissible};C\ref{cor:hardness-exact-certainty-inflation}}{CR}

\begin{definition}[Raw-vs-Certified Bit Accounting]\label{def:raw-certified-bits}
Fix a declared contract and report type $r$.
A report-bit model declares:
\begin{itemize}
\item raw report bits $B_{\mathrm{raw}}(r)$ (payload bits in the asserted report token),
\item certificate bits $B_{\mathrm{cert}}(r)$ (bits allocated to the matching evidence class).
\end{itemize}
Evidence-gated overhead and certified total are then:
\[
B_{\mathrm{over}}(r)=
\begin{cases}
B_{\mathrm{cert}}(r), & \text{if evidence for }r\text{ exists},\\
0, & \text{otherwise},
\end{cases}
\qquad
B_{\mathrm{total}}(r)=B_{\mathrm{raw}}(r)+B_{\mathrm{over}}(r).
\]
\leanmeta{\LH{IC5}, \LH{IC9}, \LH{IC12}.}
\end{definition}

\begin{proposition}[Raw/Certified Bit Split]\label{prop:raw-certified-bit-split}
For every typed report $r$:
\[
B_{\mathrm{total}}(r)=B_{\mathrm{raw}}(r)+B_{\mathrm{over}}(r),
\]
with
\[
\neg\exists\ \text{evidence for }r \Rightarrow B_{\mathrm{over}}(r)=0,\qquad
\exists\ \text{evidence for }r \Rightarrow B_{\mathrm{over}}(r)=B_{\mathrm{cert}}(r).
\]
Hence $B_{\mathrm{raw}}(r)\le B_{\mathrm{total}}(r)$ always, and strict inequality holds when evidence exists with positive certificate-bit allocation.
\leanmeta{\LH{CC11}, \LHrng{IC}{10}{11}, \LHrng{IC}{13}{16}, \LHrng{IC}{18}{19}.}
\end{proposition}

\begin{theorem}[Exact Report: Certified-Bit Gap iff Admissibility]\label{thm:exact-certified-gap-iff-admissible}
Under a report-bit model with positive exact certificate-bit allocation:
\[
\text{ClaimAdmissible}(\mathsf{EXACT})
\iff
B_{\mathrm{raw}}(\mathsf{EXACT})<B_{\mathrm{total}}(\mathsf{EXACT}).
\]
Equivalently:
\[
B_{\mathrm{total}}(\mathsf{EXACT})=B_{\mathrm{raw}}(\mathsf{EXACT})
\iff
\text{ExactCertaintyInflation}.
\]
\leanmeta{\LHrng{CC}{17}{18}, \LH{CC20}, \LH{IC31}.}
\end{theorem}

\begin{corollary}[Finite Budget Eventually Blocks Exact Admissibility]\label{cor:finite-budget-no-exact-admissibility}
\claimstamp{C\ref{cor:finite-budget-no-exact-admissibility}}{H=exp-lb-conditional}
Fix a declared physical computer model with finite budget and positive per-bit cost, a declared operation-requirement lower-bound family, and a declared typed-report bit model.
If:
\begin{enumerate}
\item exact admissibility implies certified-bit footprint at least the declared operation requirement at size \(n\), and
\item exact admissibility implies physical feasibility of that certified footprint under the declared budget model,
\end{enumerate}
then there exists \(n_0\) such that for all \(n \ge n_0\), exact admissibility fails.
The same conclusion holds in the canonical exponential-requirement specialization.
\leanmeta{\LHrng{PBC}{7}{8}.}
\end{corollary}

\begin{corollary}[Hardness-Blocked Exact Reports Are Raw-Only]\label{cor:hardness-raw-only-exact}
If exact reporting is inadmissible on the declared contract, then exact accounting collapses to raw-only:
\[
B_{\mathrm{total}}(\mathsf{EXACT})=B_{\mathrm{raw}}(\mathsf{EXACT}).
\]
\leanmeta{\LH{CC21}, \LH{IC23}.}
\end{corollary}
\claimstamp{D\ref{def:raw-certified-bits};T\ref{thm:exact-certified-gap-iff-admissible};P\ref{prop:raw-certified-bit-split};C\ref{cor:hardness-raw-only-exact}}{AR}

\begin{corollary}[No Uncertified Exact-Exhaustion Claim]\label{cor:no-uncertified-exact-claim}
If exact competence does not hold on the declared regime/objective, then an $\mathsf{EXACT}$ report is inadmissible.
\leanmeta{\LH{CC41}.}
\end{corollary}
\claimstamp{P\ref{prop:typed-claim-admissibility};C\ref{cor:no-uncertified-exact-claim}}{AR}

\begin{proposition}[Declared-Contract Separation: Selection vs. Validity]\label{prop:declared-contract-selection-validity}
For solver-side reporting, there are two formally distinct layers:
\begin{itemize}
\item \textbf{Declared-contract selection layer:} choosing objective family $(\mathcal{R}_\varepsilon)_{\varepsilon\ge 0}$, regime $\Gamma$, and assumption profile;
\item \textbf{Validity layer:} once declared, admissibility of $\mathsf{EXACT}/\mathsf{EPSILON}(\varepsilon)/\mathsf{ABSTAIN}$ reports is fixed by certificate-typed rules.
\end{itemize}
Thus the framework treats contract selection as an external declaration choice, and then enforces mechanically checked admissibility within that declared contract.
\leanmeta{\LH{CC22}, \LH{CC39}, \LH{CC75}.}
\end{proposition}

\begin{proof}
The validity layer is Proposition~\ref{prop:typed-claim-admissibility}: report admissibility is exactly tied to the matching competence certificate type.
For exact physical claims outside carve-outs, Proposition~\ref{prop:outside-excuses-explicit-assumptions} and Theorem~\ref{thm:claim-integrity-meta} require explicit assumptions, and Corollary~\ref{cor:outside-excuses-no-exact-report} blocks uncertified exact reports under declared non-collapse assumptions.
Therefore, once the contract is declared, admissibility is mechanically determined by the typed rules and attached assumption profile.
\end{proof}
\claimstamp{T\ref{thm:claim-integrity-meta};P\ref{prop:typed-claim-admissibility},\ref{prop:outside-excuses-explicit-assumptions};C\ref{cor:outside-excuses-no-exact-report}}{DC}

\begin{definition}[Exact-Claim Excuse Carve-Outs]\label{def:exact-claim-excuses}
For exact physical claims in this framework, we declare four explicit carve-outs:
\begin{enumerate}
\item trivial scope (empty in-scope family),
\item tractable class (exact competence available in the active regime),
\item stronger oracle/regime shift (the claim is made under strictly stronger access assumptions),
\item unbounded resources (budget bounds intentionally removed).
\end{enumerate}
\end{definition}

\begin{definition}[Explicit Hardness-Assumption Profile]\label{def:explicit-hardness-profile}
An explicit hardness-assumption profile for declared class/regime $(\mathcal{R},\Gamma)$ consists of:
\begin{itemize}
\item a non-collapse hypothesis,
\item a collapse implication from universal exact competence on $(\mathcal{R},\Gamma)$.
\end{itemize}
\end{definition}

\begin{definition}[Well-Typed Exact Physical Claim Policy]\label{def:exact-claim-welltyped}
An exact physical claim is \emph{well-typed} only if either:
\begin{itemize}
\item at least one carve-out from Definition~\ref{def:exact-claim-excuses} is explicitly declared, or
\item an explicit hardness-assumption profile (Definition~\ref{def:explicit-hardness-profile}) is explicitly declared.
\end{itemize}
\end{definition}

\begin{proposition}[Outside Carve-Outs, Explicit Assumptions Are Required]\label{prop:outside-excuses-explicit-assumptions}
If an exact physical claim is well-typed (Definition~\ref{def:exact-claim-welltyped}) and no carve-out from Definition~\ref{def:exact-claim-excuses} applies, then an explicit hardness-assumption profile must be present.
\leanmeta{\LH{CC22}.}
\end{proposition}

\begin{theorem}[Claim-Integrity Meta-Theorem]\label{thm:claim-integrity-meta}
For declared class/regime objects, every admissible exact-claim policy outside the carve-outs requires explicit assumptions:
\[
\begin{aligned}
&\neg\,\text{Excused (Definition~\ref{def:exact-claim-excuses})}
\ \wedge\
\text{Well-Typed Exact Claim (Definition~\ref{def:exact-claim-welltyped})} \\
&\Rightarrow\
\text{Has Explicit Hardness Profile (Definition~\ref{def:explicit-hardness-profile})}.
\end{aligned}
\]
\leanmeta{\LH{CC22}.}
\end{theorem}

\begin{proof}
Immediate from Proposition~\ref{prop:outside-excuses-explicit-assumptions}; this is the universally quantified theorem-level restatement over the typed claim/regime objects.
\end{proof}

\begin{corollary}[Outside Carve-Outs, Declared Hardness Blocks Exact Reports]\label{cor:outside-excuses-no-exact-report}
If no carve-out applies and the declared hardness-assumption profile holds, then $\mathsf{EXACT}$ reports are inadmissible for every solver on the declared class/regime.
\leanmeta{\LH{CC39}.}
\end{corollary}
\claimstamp{T\ref{thm:claim-integrity-meta};P\ref{prop:outside-excuses-explicit-assumptions};C\ref{cor:outside-excuses-no-exact-report}}{CR}

\begin{proposition}[Integrity--Competence Separation]\label{prop:integrity-competence-separation}
Integrity and competence are distinct predicates: integrity constrains asserted outputs, while competence adds non-abstaining coverage under resource bounds.
\leanmeta{\LH{IC18}, \LH{IC25}.}
\end{proposition}

\begin{proof}
Take the always-abstain solver $Q_\bot(x)=\mathsf{ABSTAIN}$ with any polynomial-time checker $V$. Definition~\ref{def:solver-integrity} holds vacuously, so $(Q_\bot,V)$ is integrity-preserving, but it fails Definition~\ref{def:competence-regime} whenever $\mathcal{X}_\Gamma\neq\emptyset$ because coverage fails. Hence integrity does not imply competence. The converse is immediate because competence includes integrity as a conjunct.
\end{proof}

\begin{definition}[Attempted Competence Failure]\label{def:attempted-competence-failure}
Fix an exact objective under regime $\Gamma$. A solver state is an \emph{attempted competence failure} if:
\begin{itemize}
\item integrity holds (Definition~\ref{def:solver-integrity}),
\item exact competence was actually attempted for the active scope/objective,
\item competence on $\Gamma$ fails for that exact objective (Definition~\ref{def:competence-regime}).
\end{itemize}
\end{definition}

\begin{proposition}[Attempted-Competence Rationality Matrix]\label{prop:attempted-competence-matrix}
Let $I,A,C\in\{0,1\}$ denote integrity, attempted exact competence, and competence available in the active regime. Policy verdict for persistent over-specification is:
\begin{itemize}
\item if $I=0$: inadmissible,
\item if $I=1$ and $(A,C)=(1,0)$: conditionally rational,
\item if $I=1$ and $(A,C)\in\{(0,0),(0,1),(1,1)\}$: irrational for the same verified-cost objective.
\end{itemize}
Hence, in the integrity-preserving plane ($I=1$), exactly one cell is rational and three are irrational.
\leanmeta{\LHrng{IC}{6}{7}, \LH{IC27}}
\end{proposition}

This separation plus Proposition~\ref{prop:attempted-competence-matrix} is load-bearing for the regime-conditional trilemma used later: if exact competence is blocked by hardness in a declared regime after an attempted exact procedure, integrity forces one of three responses: abstain, weaken guarantees, or change regime assumptions.\claimstamp{P\ref{prop:attempted-competence-matrix}}{AR}

\paragraph{Mechanized status.}
This separation is machine-checked in \texttt{DecisionQuotient/IntegrityCompetence.lean} via:
\LH{IC18} and \LH{IC25}; the attempted-competence matrix is mechanized via
\LH{IC27}, \LH{IC7}, and \LH{IC6}.

\begin{proposition}[Integrity-Resource Bound]\label{prop:integrity-resource-bound}
Let $\Gamma$ be a declared regime whose in-scope family includes all instances of SUFFICIENCY-CHECK and whose declared resource bound is polynomial in the encoding length. If $\Pclass \neq \coNP$, then no certifying solver is simultaneously integrity-preserving and competent on $\Gamma$ for exact sufficiency.
Equivalently, for every integrity-preserving solver, at least one of the competence conjuncts must fail on $\Gamma$: either full non-abstaining coverage fails or the declared budget bound fails.
\leanmeta{\LH{IC24}, \LH{IC26}, \LH{CC30}.}
\end{proposition}

\begin{proof}
By Definition~\ref{def:competence-regime}, competence on $\Gamma$ requires integrity, full in-scope coverage, and budget compliance. Under the coNP-hardness transfer for SUFFICIENCY-CHECK (Theorem~\ref{thm:sufficiency-conp}), universal competence on this scope would induce a polynomial-time collapse to $\Pclass=\coNP$. Therefore, under $\Pclass \neq \coNP$, the full competence predicate cannot hold.
Since integrity alone is satisfiable (Proposition~\ref{prop:integrity-competence-separation}, via always-abstain), the obstruction is exactly competence coverage/budget under the declared regime.
\claimstamp{T\ref{thm:sufficiency-conp};P\ref{prop:integrity-competence-separation},\ref{prop:integrity-resource-bound}}{AR}
\end{proof}

\begin{proposition}[Declared-Physics No-Universal-Exact-Certifier Schema]\label{prop:physics-no-universal-exact}
Fix any declared physical/task class representable by relation $\mathcal{R}$ and regime $\Gamma$ in the certifying-solver model. If
\[
\big(\exists Q,\ \mathrm{CompetentOn}(\mathcal{R},\Gamma,Q)\big)\ \Rightarrow\ (\Pclass=\coNP),
\]
then under $\Pclass\neq\coNP$ there is no universally exact-competent solver for that declared class.
\leanmeta{\LH{CC15}.}
\end{proposition}

\begin{corollary}[No Universal Exact-Claim Admissibility in Hardness-Blocked Classes]\label{cor:physics-no-universal-exact-claim}
Under the same premises as Proposition~\ref{prop:physics-no-universal-exact}, and under the typed-claim discipline of Definition~\ref{def:typed-claim-report}, an $\mathsf{EXACT}$ report is inadmissible for every solver on that declared class.
Hence admissible reporting must use $\mathsf{ABSTAIN}$ or explicitly weakened guarantees.
\leanmeta{\LH{CC38}.}
\end{corollary}
\claimstamp{P\ref{prop:physics-no-universal-exact};C\ref{cor:physics-no-universal-exact-claim}}{CR}

Informally: honesty limits what you can claim, capability limits what you can solve, and hard regimes may require abstaining or using weaker exact claims.


% Include the hardness proofs (already developed)
\input{hardness_proofs.tex}
